\documentclass[a4paper,11pt]{article}

\usepackage[T1]{fontenc}
\usepackage[utf8]{inputenc}
\usepackage[slovene]{babel}
\usepackage{lmodern}
\usepackage{amsmath,amssymb,amsthm}
\usepackage{graphicx}
\usepackage[margin=2.5cm]{geometry}
\usepackage[hidelinks]{hyperref}

\newtheorem{izrek}{Izrek}[section]
\theoremstyle{definition}
\newtheorem{definicija}[izrek]{Definicija}
\theoremstyle{remark}
\newtheorem{opomba}[izrek]{Opomba}

\title{Aproksimacija krožnic z Bézierjevimi krivuljami\break
       z zvezno ukrivljenostjo}
\author{Ime in priimek\\
  \small Seminar pri predmetu Računalniško podprto geometrijsko oblikovanje}
\date{\today}

\begin{document}
\maketitle

\begin{abstract}
% Na kratko predstavite problem, uporabljeno metodo in glavne rezultate.
Članek obravnava aproksimacijo krožnih lokov z Bézierjevimi krivuljami
s poudarkom na zveznosti ukrivljenosti. Predstavljeni bodo geometrijski
pogoji konstrukcije, merjenje aproksimacijske napake in numerični primeri.
\end{abstract}

\noindent\textbf{Ključne besede:} Bézierjeva krivulja, krožni lok,
aproksimacija, ukrivljenost, geometrijska zveznost.

\section{Uvod}
\label{sec:uvod}
% Opišite motivacijo, cilj seminarja in zgradbo članka.
% Dodajte popolne bibliografske podatke izhodiščnega članka v literaturo.

\section{Matematične osnove}
\label{sec:osnove}

\subsection{Bézierjeve krivulje}
Bézierjeva krivulja stopnje $n$ s kontrolnimi točkami
$\mathbf{b}_0,\ldots,\mathbf{b}_n\in\mathbb{R}^2$ je podana z
\begin{equation}
  \mathbf{b}(t)=\sum_{i=0}^{n}\mathbf{b}_i B_i^n(t),
  \qquad B_i^n(t)=\binom{n}{i}t^i(1-t)^{n-i},
  \qquad t\in[0,1].
\end{equation}
% Dopolnite lastnosti in izraze za odvode.

\subsection{Ukrivljenost in geometrijska zveznost}
Za regularno ravninsko krivuljo $\mathbf{r}(t)=(x(t),y(t))$
je predznačena ukrivljenost enaka
\begin{equation}
  \kappa(t)=
  \frac{x'(t)y''(t)-y'(t)x''(t)}
       {\bigl(x'(t)^2+y'(t)^2\bigr)^{3/2}}.
\end{equation}
% Razložite pogoje G^0, G^1 in G^2 ter razliko do parametrične zveznosti.

\section{Postavitev aproksimacijskega problema}
\label{sec:problem}
% Določite polmer, kot loka, simetrijo in robne pogoje.
Krožni lok s središčem v izhodišču in polmerom $R>0$ opišemo kot
\begin{equation}
  \mathbf{c}(\varphi)=R(\cos\varphi,\sin\varphi),
  \qquad \varphi\in[\alpha,\beta],\quad \alpha<\beta.
\end{equation}
% Natančno opredelite uporabljeno mero napake in dopustne krivulje.

\section{Konstrukcija aproksimacijske krivulje}
\label{sec:konstrukcija}
\subsection{Izbira kontrolnih točk in robni pogoji}
% Izpeljite pogoje za krajišči, tangenti in ukrivljenost.

\subsection{Določitev prostih parametrov}
% Predstavite enačbe, kriterij izbire in postopek reševanja.

\section{Analiza napake}
\label{sec:napaka}
% Izpeljite napako in obravnavajte njene ekstreme.
% Ločite radialno napako od drugih mer geometrijske razdalje.

\section{Numerični primeri}
\label{sec:primeri}
% Navedite parametre, uporabljeno programsko opremo in natančnost računanja.
% Dodajte grafe krivulj, ukrivljenosti in napake ter tabelo rezultatov.
% Primerjajte več kotov loka in rezultate komentirajte.

\section{Zaključek}
\label{sec:zakljucek}
% Povzemite ugotovitve, omejitve metode in možnosti nadaljnjega dela.

\begin{thebibliography}{9}
% Vzorec zapisa (odkomentirajte in dopolnite s preverjenimi podatki):
% \bibitem{izhodišče}
% Avtorji, \emph{Good approximation of circles by curvature-continuous
% Bézier curves}, revija, letnik, strani, leto. DOI.
\end{thebibliography}

\end{document}
